%========================================================================%
\section{Methods}\label{sec:methods}
%========================================================================%

\begin{figure*}[t]
\centering
\resizebox{0.95\textwidth}{!}{%
\begin{tikzpicture}[
  font=\sffamily\small,
  >={Stealth[length=2.2mm,width=1.8mm]},
  node distance=6mm and 8mm,
  windowbox/.style={rectangle, draw=blue!55, rounded corners=2pt,
        minimum width=18mm, minimum height=8mm,
        fill=blue!8, font=\sffamily\footnotesize, inner sep=2pt},
  stagebox/.style={rectangle, draw, rounded corners=3pt,
        minimum width=66mm, minimum height=24mm, align=center,
        font=\sffamily\small, inner sep=4pt},
  s1box/.style={stagebox, draw=green!55!black, fill=green!10},
  s2box/.style={stagebox, draw=orange!75!black, fill=orange!12},
  cachebox/.style={rectangle, draw=purple!70, rounded corners=3pt,
        minimum width=70mm, minimum height=11mm,
        fill=purple!10, font=\sffamily\small, align=center},
  quantbox/.style={rectangle, draw=black!60, dashed, rounded corners=2pt,
        minimum width=34mm, minimum height=11mm,
        fill=yellow!12, font=\sffamily\footnotesize, align=center},
  ann/.style={font=\sffamily\scriptsize, text=black!60, align=center},
  flow/.style={->, gray!65, line width=0.4pt},
  flowQ/.style={->, black!85, line width=0.7pt},
  flowS1/.style={->, green!50!black, line width=0.7pt},
  flowS2/.style={->, orange!75!black, line width=0.7pt},
]

% --- Row 1: input windows (centered above the diagram) ---
\node[windowbox] (w1) at (0,0)          {$w_1$};
\node[windowbox, right=4mm of w1] (w2)  {$w_2$};
\node[windowbox, right=4mm of w2] (wd)  {$\cdots$};
\node[windowbox, right=4mm of wd] (wT)  {$w_T$};
\node[windowbox, right=10mm of wT, draw=black!75, fill=blue!18, very thick] (wq)
        {$w_q$ (query)};

\node[ann, above=1mm of w2.north east, anchor=south]
        {context windows};
\node[ann, above=1mm of wq.north, anchor=south]
        {current question};

% --- Row 2: Stage 1 and Stage 2, side by side, centered under windows ---
\coordinate (rowcenter) at ($(w1.south)!0.5!(wq.south)$);

\node[s1box, below=14mm of rowcenter, xshift=-38mm]
  (s1) {%
    \textbf{Stage 1.} Parameter-free selection \\[1pt]
    $\mathrm{score}_i = \max(h_i,0)\,\cdot\,\max(c_i,0)$ \\[2pt]
    {\footnotesize cumulative attention $h_i$ $\times$}\\
    {\footnotesize query--key alignment $c_i$}\\[1pt]
    \textbf{top-$k$ real tokens} ($k_{\mathrm{real}}{=}284$)
  };

\node[s2box, right=12mm of s1]
  (s2) {%
    \textbf{Stage 2.} Learned summary \\[1pt]
    {16 memory tokens, EMA aggregator}\\[2pt]
    {\footnotesize hidden-state sensing across windows}\\
    {\footnotesize virtual K/V via frozen $\mathbf{W}_K, \mathbf{W}_V$}\\[1pt]
    \textbf{trained}: $\sim$14\,M params (only module)
  };

% --- Row 3: combined cache (centered under both stages) ---
\coordinate (stagemid) at ($(s1.south)!0.5!(s2.south)$);
\node[cachebox, below=12mm of stagemid]
  (cache) {%
    \textbf{Combined KV cache} ($k = 300$ tokens) $\to$ frozen LLM \\[1pt]
    {\footnotesize $25.6\times$ compression vs.\ full cache}
  };

% --- Right-side optional quantisation ---
\node[quantbox, right=10mm of cache]
  (kvq) {%
    K8V4 quantisation\\[1pt]
    {\footnotesize +$2.6\times$ compression}
  };

% --- Arrows: windows to Stage 1 (gray, context only) ---
\foreach \src in {w1,w2,wd,wT}{
  \draw[flow] (\src.south) -- ($(\src.south)+(0,-3mm)$) -| (s1.north);
}
% --- Arrows: windows to Stage 2 (gray, context only) ---
\foreach \src in {w1,w2,wd,wT}{
  \draw[flow] (\src.south) -- ($(\src.south)+(0,-3mm)$) -| (s2.north);
}
% --- Query: emphasised path into Stage 1 only ---
\draw[flowQ] (wq.south) -- ($(wq.south)+(0,-3mm)$)
   -| (s1.north east)
   node[pos=0.92, right=1mm, font=\sffamily\scriptsize, text=black!75]
   {query-aware};

% --- Stage 1 output (real tokens) into cache ---
\draw[flowS1] (s1.south) -- ($(s1.south)+(0,-4mm)$)
   -| ($(cache.north)+(-18mm,0)$)
   node[pos=0.92, left=1mm, font=\sffamily\scriptsize, text=green!40!black]
   {284 real};

% --- Stage 2 output (summary tokens) into cache ---
\draw[flowS2] (s2.south) -- ($(s2.south)+(0,-4mm)$)
   -| ($(cache.north)+(18mm,0)$)
   node[pos=0.92, right=1mm, font=\sffamily\scriptsize, text=orange!75!black]
   {16 summary};

% --- Optional quantisation arrow ---
\draw[->, dashed, black!60] (cache.east) -- (kvq.west);

\end{tikzpicture}%
}
\caption{\textbf{The AstroNet pipeline.} A frozen quantised backbone
processes $T$ context windows $w_1,\ldots,w_T$ followed by a query window
$w_q$. \textbf{Stage 1} (green) is parameter-free: every cached token
receives a score that multiplies its cumulative attention $h_i$ by its
cross-window query--key alignment $c_i$ (\cref{sec:s1}), and the top
$k_{\mathrm{real}}{=}284$ tokens are retained. \textbf{Stage 2} (orange)
is the only trained module: sixteen memory tokens aggregate hidden states
across windows via an EMA state and emit virtual key--value pairs
through dequantised copies of the backbone's own
$\mathbf{W}_K, \mathbf{W}_V$ projections (\cref{sec:s2}). The combined
$300$-token cache is consumed by the frozen backbone for generation
and composes with downstream KV-cache quantisation (dashed,
\cref{sec:quant}) for a further $\sim 2.6\times$ compression.}
\label{fig:architecture}
\end{figure*}

%-----------------------------------------------------------------------%
\subsection{Setup}\label{sec:setup}
%-----------------------------------------------------------------------%

We evaluate six frozen autoregressive transformer backbones spanning three
architecture families and a 4.5$\times$ range in parameter count: Qwen 2.5
at 7B, 14B and 32B; Llama 3.1-8B; Mistral 7B v0.3; and Mistral-Small at
24B. All six are loaded in 4-bit NF4 quantisation through
bitsandbytes\cite{dettmers2023nf4} with float16 compute, the standard
deployment configuration for single-GPU serving in this size class. The
backbones are not modified at any point: no weight is updated, no
architectural component is replaced, and the AstroNet wrapper interacts
with the model only through forward hooks on the residual stream and
post-softmax attention tensors, both of which are exposed by the
\texttt{transformers} library through the eager attention implementation
without requiring custom kernels.

Following the multi-window convention used in long-context cache
studies, each evaluation instance is presented as a sequence of $T$ fact
windows $w_1,\dots,w_T$ followed by a single query window $w_q$.
Processing a fact window $w_t$ at every layer $l$ produces a set of
key--value pairs $\{(\mathbf{k}_i^{(l)},\mathbf{v}_i^{(l)})\}_{i=1}^{|w_t|}$.
Given a total cache budget of $k$ tokens, the task is to decide which
KV entries to retain across all $T$ fact windows so that the model can
answer the query in $w_q$. Unless stated otherwise we use $k=300$ for
all selection methods; this is small relative to the $\sim$1500-token
total context of the controlled SQuAD setup and is the regime in which
selection strategies are stressed. Experiments were run on a workstation
with two NVIDIA Titan RTX 24~GiB GPUs using PyTorch 2.4.1 and
transformers 4.46.3; models from 32B upwards are split across both GPUs
with a layer-balanced device map.


%-----------------------------------------------------------------------%
\subsection{Stage 1: parameter-free multiplicative selection}\label{sec:s1}
%-----------------------------------------------------------------------%

Stage 1 scores every cached token by a product of two complementary
signals derived from the frozen model's own state and discards all but
the top-$k$. Both signals are computed at inference time from quantities
the model already produces; no parameters are introduced.

The first signal is the cumulative attention each token has received
during processing of the fact windows. For token $i$ we sum the
post-softmax attention weights it received over all queries $q$, heads
$h$, and layers $l$:
\begin{equation}
    h_i = \sum_l \sum_h \sum_q a_i^{(l,h,q)}.
    \label{eq:heuristic}
\end{equation}
This is the heavy-hitter heuristic of H${}_2$O\cite{zhang2023h2o},
which captures how strongly the model's attention engaged with each
token while reading the context but is blind to the query that will
eventually be asked.

The second signal is the cross-window relevance of each cached token to
the query. Once the query window has been processed we extract the
query projections $\mathbf{q}_{\mathrm{query}}^{(l,h)}$ at a small set
$\mathcal{L}$ of four evenly spaced layers and score each cached token
against them:
\begin{equation}
    c_i = \sum_{l \in \mathcal{L}} \sum_h
          \mathbf{q}_{\mathrm{query}}^{(l,h)} \cdot \mathbf{k}_i^{(l,h)}
          \big/ \sqrt{d_h}.
    \label{eq:cross}
\end{equation}
The set $\mathcal{L}$ contains the same layers at which Stage 2 (below)
injects, chosen to balance early and late representations; the
sensitivity to this choice is reported in
\cref{sec:ablations}. Computing $c_i$ requires one forward pass on
$w_q$ plus $|\mathcal{L}|$ matrix multiplications, an overhead
substantially smaller than a single additional forward pass.

The two signals are combined multiplicatively rather than additively:
\begin{equation}
    \mathrm{score}_i = \max(h_i, 0)\cdot\max(c_i, 0),
    \label{eq:mult}
\end{equation}
and the $k$ tokens of highest score are retained, ordered by their
original position. The product enforces conjunctive selection: a token
is kept only if it was salient during processing \emph{and} relevant
to the query.

We chose a multiplicative rather than additive combination because the
two signals capture orthogonal information that a sum would mix: the
heuristic $h_i$ measures \emph{how much the model engaged with this token
during processing} (regardless of any future query), while $c_i$
measures \emph{how relevant this token is to the current query}. An
additive combination $h_i + c_i$ treats them as comparable contributions
to a single score; a multiplicative combination $h_i \cdot c_i$ requires
both signals to be jointly large, suppressing tokens that look
locally salient but are unrelated to the question. Empirically this
matters: the cumulative-attention-only H${}_2$O baseline (which is
$h_i$ alone) falls 30--47 percentage points below the multiplicative
score on every backbone we evaluate, and a matched-budget additive
ablation underperforms the multiplicative form by 8--12 points
(\cref{sec:ablations}).

We note that Stage 1 is query-aware: it consults $w_q$ before deciding
which cached entries to keep, in contrast to H${}_2$O, SnapKV and
StreamingLLM, which select using only processing-time attention. This
brings the family of comparisons in our experiments into two groups,
query-unaware and query-aware (which includes RAG); we report both.


%-----------------------------------------------------------------------%
\subsection{Stage 2: learned cross-window summary tokens}\label{sec:s2}
%-----------------------------------------------------------------------%

Stage 1 operates one token at a time and cannot represent information
whose support spans multiple windows. Stage 2 addresses this with a
small learned module that maintains a persistent state $\mathbf{g}$
aggregating window-level summaries through an exponential moving
average, and emits a fixed number of virtual KV pairs from $\mathbf{g}$
that are prepended to the Stage-1 selection. The total cache budget is
held fixed: with $K=16$ summary tokens, Stage 1 retains $k-K=284$ real
tokens, so any improvement of the hybrid over Stage 1 alone is at the
same end-to-end memory cost.

After processing each fact window the module senses the model's
hidden states $\mathbf{H}\in\mathbb{R}^{S\times d}$ at a designated
sensing layer $l_s$ through cross-attention with $K=16$ learnable query
vectors $\mathbf{Q}_{\mathrm{mem}}\in\mathbb{R}^{K\times d_a}$:
\begin{equation}
    \mathbf{m} = \mathrm{softmax}\!\Bigl(
        \mathbf{Q}_{\mathrm{mem}}\cdot\mathrm{LN}(\mathbf{H})^\top
        \big/\sqrt{d_a}\Bigr)\cdot\mathrm{LN}(\mathbf{H}),
    \label{eq:sense}
\end{equation}
where $\mathrm{LN}$ is layer normalisation, applied for compatibility
with the wide dynamic range of pre-trained hidden states. The
result $\mathbf{m}\in\mathbb{R}^{K\times d}$ is a compressed, content-
aware summary of the window.

A persistent state $\mathbf{g}\in\mathbb{R}^{K\times d}$ aggregates
window-level summaries through an exponential moving average with
learnable rate $\alpha\in(0,1)$:
\begin{equation}
    \mathbf{g}(t{+}1) = (1-\alpha)\cdot\mathbf{g}(t) + \alpha\cdot\mathbf{m}(t),
    \label{eq:ema}
\end{equation}
initialised at $\mathbf{g}(0)=\mathbf{0}$ and $\alpha_0=0.38$. We
report in \cref{sec:ablations} that simpler aggregators (mean
pooling across windows, retaining only the last window) perform
comparably on five of six backbones, indicating that the specific EMA
form is not load-bearing.

The state $\mathbf{g}$ does not enter the attention computation
directly. At inference, $\mathbf{g}$ is first transformed by a
per-layer low-rank bottleneck
$\mathbb{R}^d\to\mathbb{R}^{256}\to\mathbb{R}^d$, with no
non-linearity, to produce virtual hidden states
$\tilde{\mathbf{h}}^{(l)}\in\mathbb{R}^{K\times d}$ at every layer
$l$. The virtual hidden states are then passed through dequantised
copies of the frozen model's own key and value projection matrices,
\begin{equation}
    \tilde{\mathbf{K}}^{(l)} = \tilde{\mathbf{h}}^{(l)} \mathbf{W}_K^{(l)},\quad
    \tilde{\mathbf{V}}^{(l)} = \tilde{\mathbf{h}}^{(l)} \mathbf{W}_V^{(l)},
    \label{eq:virtkv}
\end{equation}
producing $K$ virtual KV pairs at every layer. Because
$\mathbf{W}_K^{(l)}$ and $\mathbf{W}_V^{(l)}$ are the model's own
projection matrices (dequantised at the wrapper, but functionally
identical to those used to produce real KV entries), the virtual keys
and values inhabit the same representation space as real entries and
the attention computation treats them identically. The virtual pairs
are prepended to the Stage-1 selection before generation begins.

The only trainable parameters in Stage 2 are the cross-attention
queries and projections in Eq.\,\ref{eq:sense}, the EMA rate $\alpha$,
and the per-layer low-rank bottleneck weights. This totals 10--14M
parameters per backbone depending on the hidden dimension, less than
$0.2\%$ of any of the six backbones we evaluate. A soft regulariser
constrains the norm of $\tilde{\mathbf{h}}^{(l)}$ to track the
embedding-layer norm, which we found necessary for stable training
across the range of hidden dimensions in our model set; the inject
scale is held at $0.03$ of the residual norm at every layer.


%-----------------------------------------------------------------------%
\subsection{Quantisation: per-head Lloyd--Max}\label{sec:quant}
%-----------------------------------------------------------------------%

For additional cache compression we adapt the Lloyd--Max optimal scalar
quantisation of TurboQuant\cite{zandieh2026turboquant}. TurboQuant has
two stages: an MSE-optimal stage that applies a random orthogonal
rotation followed by per-coordinate Lloyd--Max scalar quantisation, and
a residual-correction stage that uses quantised
Johnson--Lindenstrauss projections to recover unbiased inner-product
estimates. The residual-correction stage requires bypassing the
attention multiply with a JL-based estimator, which is incompatible
with frozen-attention kernels; we use only the MSE stage, which
produces dequantised key and value tensors that the standard attention
path consumes without modification.

The random rotation that TurboQuant applies in the MSE stage destroys
the paired-coordinate structure that rotary position
encodings\cite{su2024rope} rely on: RoPE rotates pairs of channels by
position-dependent angles, and an arbitrary orthogonal rotation across
all channels scrambles this pairing. In our hands, vanilla
TurboQuant at 4 bits reduces accuracy on Qwen 7B to zero on the
multi-window SQuAD task. The same failure mode is structurally
expected on every RoPE-using backbone we evaluate.

Our adapted variant removes the rotation entirely. We instead
normalise each attention head's key and value coordinates
independently to zero mean and unit variance using head-level running
statistics, apply Lloyd--Max scalar quantisation with a codebook tuned
to the standard normal distribution, and denormalise. We use 8-bit
keys and 4-bit values (K8V4): keys participate in inner products
$\mathbf{q}\cdot\mathbf{k}$ that are sensitive to coordinate noise,
while values enter only through attention-weighted averages that
tolerate coarser representation. Per token retained, the K8V4
encoding uses $1.5\times d_h$ bytes versus $4\times d_h$ for FP16,
a $\sim$2.6$\times$ reduction; combined with the
$10\times$ selection compression for a 20-window context relative
to the full cache, the total is $\sim$26$\times$ at FP16 and
$\sim$68$\times$ with K8V4 (Table~\ref{tab:compression}). The
per-head normalisation step is
selection-method-agnostic: it can be applied to the KV tensors
retained by H${}_2$O, SnapKV, StreamingLLM or any other selection
scheme on a RoPE backbone, and the fix is a general contribution
usable beyond AstroNet.


%-----------------------------------------------------------------------%
\subsection{Training}\label{sec:training}
%-----------------------------------------------------------------------%

Only the Stage-2 module is trained; the backbone, its quantised
weights, and Stage 1 are frozen throughout. Two training protocols are
used. The first uses 5{,}000 SQuAD training
examples\cite{rajpurkar2016squad} arranged as five-window instances
(four 256-token fact windows plus a query window), trained for one
to three epochs of next-token cross-entropy on the answer span with
AdamW at a learning rate of $5\times10^{-5}$, batch size one and
four-step gradient accumulation. This protocol produced the Qwen 7B,
Qwen 14B, Qwen 32B, Llama 3.1-8B and Mistral 7B checkpoints used
throughout the paper.

For Mistral-Small 24B this protocol produced a Stage-2 module that
slightly regressed accuracy relative to Stage 1 alone at evaluation
context lengths longer than the training distribution. We therefore
adopted an extended protocol for the two Mistral backbones: a
5{,}000-sample corpus drawn in equal proportions from SQuAD and
HotpotQA, presented as ten-window instances rather than five, with the
rest of the training configuration unchanged. The extended protocol
resolves the 24B regression and is documented in the Discussion as a
transparent finding about training-distribution coverage for larger
backbones rather than as an architecture-specific modification; the
smaller backbones reuse their original checkpoints. The training corpus
saturates at 5{,}000 samples: increasing to 10{,}000 or 20{,}000
samples does not improve held-out accuracy.


%-----------------------------------------------------------------------%
\subsection{Baselines}\label{sec:baselines}
%-----------------------------------------------------------------------%

We compare against seven baselines covering the two prior families
identified in the Introduction. The cache-eviction family comprises
five methods. H${}_2$O\cite{zhang2023h2o} scores tokens by cumulative
attention received during processing and retains a heavy-hitter
budget together with a window of the most recent tokens. SnapKV\cite{li2024snapkv}
uses an observation window at the end of the prompt to identify
positions whose post-softmax attention is locally concentrated,
applying average-pool smoothing and a recent-token bias before
selection. StreamingLLM\cite{xiao2024streamingllm} retains a small
number of initial attention sinks together with a sliding window of
the most recent tokens. PyramidKV\cite{cai2024pyramidkv} allocates a
layer-adaptive budget that decreases with depth, motivated by an
observed concentration of useful attention in shallow layers.
KIVI\cite{liu2024kivi} compresses the retained cache through
asymmetric integer quantisation, with per-channel quantisation for
keys and per-token quantisation for values; we evaluate KIVI at 4
bits (K4V4) as the operating point matched in cache-byte cost to our
K8V4 stack.

Two retrieval baselines round out the comparison. RAG~$k{=}1$ and
RAG~$k{=}3$ encode each fact window with the
\texttt{all-MiniLM-L6-v2} sentence-transformer (22M parameters,
fully external to the backbone), score windows against an encoding of
the query, and prepend the top-$k$ retrieved windows to the query
before generation. In our five-window setup, RAG~$k{=}1$ recovers the
single best-matching window out of four fact windows and RAG~$k{=}3$
recovers three of four; RAG~$k{=}3$ is intentionally generous and
serves as an upper bound for in-context retrieval at this scale.

H${}_2$O and SnapKV in their original formulations select tokens
per attention head. We re-implement them with a global-pooling
selection rule that scores each token by averaging over heads before
ranking. This decision is forced by the small per-window budget of
our setup: at $k=300$ tokens distributed across all layers and
attention heads, per-head selection scatters the cache so thinly that
both methods collapse to $\sim$8.5\% accuracy on Qwen 7B, well
below their global-pooled performance. Global pooling is the variant
used throughout the main tables; the per-head ablation is reported in
the Supplementary Information as a methodological transparency point.


%-----------------------------------------------------------------------%
\subsection{Evaluation protocols}\label{sec:eval}
%-----------------------------------------------------------------------%

Five evaluation regimes are used in the paper. The controlled
multi-window QA protocol on SQuAD\cite{rajpurkar2016squad} is the
position-robust variant: each instance is presented in four
configurations differing only in which of the four fact windows
contains the answer, repeated across five random seeds, with 100
samples per (seed, position) cell. Reported accuracy on this protocol
is averaged over all twenty (seed, position) pairs, with confidence
intervals computed across seeds at fixed position. LongBench\cite{bai2024longbench}
is used as a transfer benchmark on the HotpotQA and MultiFieldQA
subsets: 100 instances per task, documents chunked into 384-token
windows and processed sequentially through the same pipeline used at
SQuAD evaluation, with F1 computed against gold answers after first-
newline truncation. Needle-in-a-Haystack is run at five depth
positions of the answer-bearing fact among $n\in\{5,10,20\}$ distractor
windows with 20 trials per depth. The long-context generalisation
benchmark uses the RULER\cite{hsieh2024ruler} multi-hop subset at
$n=22$, $44$ and $85$ windows, corresponding approximately to 8k, 16k
and 32k input tokens; the Stage-2 training distribution covers up to
$n=10$, so the RULER points at 16k and 32k probe genuine
out-of-distribution context lengths. Finally, latency is measured as
forward-pass wall-clock time at fixed cache budget on Qwen 7B and
Llama 8B over twenty samples; this isolates the structural cost of
Stage 1's dependence on post-softmax attention scores, which precludes
memory-efficient attention kernels.
